\chapter{Standard Model Effective Field Theory (SMEFT)}\label{ch:smeft}
For the energy regime above the electroweak scale, the full Standard Model Effective Field Theory (SMEFT) is the appropriate effective description. Since we are exploring the imprints of heavy new physics on the phenomenology of muon decay we first have to identify the set of relevant SMEFT operators -- those operators which result in their low-energy limit in the LEFT operators identified in the previous chapter. For this purpose we are using the tree-level matching results of \cite{jenkins2018lowOM} and \parencite{liao2020extending}. However, while the matching is performed at the electroweak scale, the investigation of potential heavy new fields possibly shielded by these effective operators naturally takes place at the energy scale of new physics -- roughly at the order of a fex $\text{TeV}$, as we know from Tab.~\ref{tab:leftbounds}. So before we can \enquote{open up} the SMEFT operators we first have to \texttt{run} them up to that energy scale. To accomplish this task we are using the python packages \texttt{wilson} \cite{aebischer2018python} and the \texttt{D7RGESolver} \parencite{liao2025rge}. This will allow us to present the full set of SMEFT operators at the new physics scale which in their low-energy limit affect muon decay.\par
The content of this chapter is ordered in the following way: First, we briefly summarize the structure of SMEFT and its operators at dimensions five, six and seven. Subsequently, we cite the tree-level matching conditions for the relevant SMEFT operators at dimension six and seven to the LEFT operators which perturb muon decay. Finally, the identified SMEFT operators are run up to the energy scale of new physics where, in the next chapter, we can open them up to explore the various heavy new fields possibly veiled by them.
\section{The SMEFT Lagrangian}
The SMEFT extends the SM Lagrangian by embedding it as the leading term of an expansion in terms of powers of a heavy scale $\Lambda$: At each canonical mass dimension $D$ (in our convention the $\Lambda$s are absorbed by the Wilson coefficients, following \cite{jenkins2013rge1}) we add all local Lorentz invariant operators built out of the SM field content:
\begin{equation}
    \mathcal{L}=\mathcal{L}_\text{SM}+\sum_{D\ge5}\mathcal{L}^D
\end{equation}
Unlike the HEFT, the SMEFT assumes that electroweak symmetry is linearly realized and that the physical Higgs belongs to the $SU(2)_L$ doublet $H$. The higher dimensional interactions are therefore constructed from complete SM muliplets and respect the full SM gauge group $SU(3)_C\times SU(2)_L\times U(1)_Y$.\par 
At dimension five, there is only one operator (class) we can add -- the Weinberg neutrino mass operator \cite{weinberg1967model}, which combines two Higgs doublets with two lepton doublets:
\begin{equation}
    \mathcal{L}^5=C_5^{rs}\epsilon ^{ik}\epsilon^{kl}(l^T_{ir}Cl_{ks})H_jH_l+\text{h.c.}
\end{equation}
Here, $r,s$ are flavor indices, $i,j,k,l$ weak-eigenstate indices. For three generations, $C_5$ has 6 complex entries. After spontaneous symmetry breaking, this operator produces a Majorana mass term for the neutrino: $M_\nu=-v^2C_5$.\par 
For dimension six, Buchmüller and Wyler counted 80 operators in 1986 \cite{buchmuller1986effective}. However, many of these operators turn out to be actually equivalent and refer to the same operator structure because they can be transformed into each other. Fierz transformations, for example -- one of which we used in Eq.~\ref{eq:fierzingVLR2SRR} to \enquote{fierz} $L_{\nu e}^{V,LR}$ into $g^S_{RR}$ -- account for many subtle equivalences of scalar, vector and tensor-like structures. Moreover, the SM equations of motion and field redefinitions also allow for transforming many operator structure into each other. An updated classification was therefore published in 2010 which is now mostly referred to as the Warsaw basis \cite{grzadkowski2010dimension} and which counts 63 independent operators, 59 of which conserve baryon number. Of those 59 operators, 42 are self-hermitian and 17 not -- which is why Jenkins, Manohar \& Stoffer count $42+2(17)=76$ baryon-number conserving operators in \cite{jenkins2018lowOM} and $4+4=8$ baryon-number violating operators. At dimension seven, there are 18 independent operators \cite{liao2020extending} -- originally Lehman~\cite{lehman2014extending} presented 20 operators, but 2 pairs of $\psi^4D$ operators were shown to refer to the same structure. In Tab.~\ref{tab:SMEFTnumOP} we collect the number of operators at dimension five, six and seven, following the counting in \cite{jenkins2018lowOM} and \cite{liao2020extending}. We cite the full bases for dimensions six and seven in the appendix in Ch.~\ref{sec:base}.
\begin{table}[t]
    \centering
    \begin{tabular}{llll} \toprule
             D&$(\Delta B,\Delta L)$&$n_g=1$&$n_g=3$\\\midrule
             5&$(0,+2)$ + h.c.&$1+1$&6 + 6\\
             6&$(0,0)$&76&2499\\
             6&$(1,1)$ + h.c.&4 + 4&273 + 273\\
             7&$(0,+2)$ + h.c.&12 + 12&474 + 474\\
             7&$(+1,-1)$ + h.c.&6 + 6&297 + 297\\\bottomrule
    \end{tabular}
    \caption{Number of SMEFT operators at dimensions five, six and seven from \cite{jenkins2018lowOM} and \cite{liao2020extending} ($n_g=$ number of generations)}
    \label{tab:SMEFTnumOP}
\end{table}
\section{SMEFT in the broken phase}
In the spontaneously broken SMEFT effective Higgs operators alter the VEV, the couplings and the weak mixing angle -- for this discussion we are following the section of the same title in \cite{jenkins2018lowOM}. Essential for muon decay, the gauge currents themselves receive operator-dependent corrections. For example $\mathcal{W}^+_\mu$ couples to $(\bar\nu_{Lp}\gamma^\mu e_{Lr})$ via a modified coupling $[W_l]_{pr}$:
\begin{equation}
    \Delta\mathcal{L}_{Wl\nu}=-\frac{\bar g_2}{\sqrt{2}}\mathcal{W}_\mu^+\bar\nu_{Lp}\gamma^\mu[W_l]_{pr}e_{Lr}+\text{h.c.}
\end{equation}
The modified coupling $[W_l]_{pr}$ as well as $[Z_\psi]_{pr}$ which couples $\mathcal{Z}_\mu$ to $(\bar\psi_p\gamma^\mu\psi_r)$ whereby $\psi=\nu_L,e_L,e_R$ in the spontaneously broken SMEFT receive corrections from $\mathcal{O}^{(3)}_{Hl}$ as well as from the modified VEV:
\begin{align}
    [W_l]_{pr} &= \left[\delta_{pr} + v_T^2  C^{(3),pr}_{Hl} \right]\\[5pt]
    [Z_{\nu_L}]_{pr} &= \left[\delta_{pr}\left(\frac 12\right) - \frac12 v_T^2
    C^{(1),pr}_{Hl} + \frac12 v_T^2  C^{(3),pr}_{Hl} \right]\\[5pt]
    [Z_{e_L}]_{pr} &= \left[\delta_{pr}\left(-\frac 12+\bar s^2 \right) - \frac12 v_T^2  C^{(1),pr}_{Hl} - \frac12 v_T^2  C^{(3),pr}_{Hl} \right]\\[5pt]
    [Z_{e_R}]_{pr} &= \left[\delta_{pr}\bar s^2  - \frac12 v_T^2  C^{pr}_{He}\right]
\end{align}
The VEV and the gauge coupling change in the following way, which also affects the weak mixing angle:
\begin{align}
    v_T = \left( 1+ \frac{3 L_H v^2}{8 \lambda} \right) v\qquad \bar g_1= g_1 \left(1 + L_{HB} \, v_T^2 \right)\qquad \bar g_2= g_2 \left(1 + L_{HW} \, v_T^2 \right)\label{eq:vtgaugecorr}\\[5pt]
    \bar s =\sin \theta =  \frac{\bar g_1}{\sqrt{ \bar g_1^2 + \bar g_2^2}} \left[ 1 + \frac{\epsilon}{2}\frac{\bar g_2}{\bar g_1} \left( \frac{\bar g_2^2 - \bar g_1^2}{\bar g_1^2 + \bar g_2^2} \right)   \right]\hspace{2.2cm}
\end{align}
And the effective couplings $\bar e$ and $\bar g_Z$ as well:
\begin{align}
\bar e &= \bar g_2 \, \sin \bar \theta - \frac{1}{2} \, \cos \bar \theta  \, \bar g_2 \, v_T^2 \, C_{HWB}, \nonumber\\[5pt]
\bar g_Z &= \frac{\bar e}{\sin \bar\theta \cos \bar\theta}  \left[1 +  \frac{\bar g_1^2+\bar g_2^2}{2 \bar g_1\bar g_2} v_T^2C_{HWB}\right]
\end{align}
Equipped with these corrections we can now take a look on the matching conditions of the SMEFT onto the LEFT.
\section{Matching conditions}\label{sec:matchingLEFTSMEFT}
From the LEFT paper by Jenkins, Manohar and Stoffer \cite{jenkins2018lowOM} on operators and matching at dimension six we collect the following matching conditions for the LNC operators $\mathcal{O}^{V,LL}_{\nu e}$ and $\mathcal{O}^{V,LR}_{\nu e}$:
\begin{align}
    L^{V,LL}_{\nu e,prst}&= C_{ll}^{prst}+C_{ll}^{stpr}-\tfrac{\bar g_2^2}{2M_W^2}[W_l]_{pt}[W_l]^*_{rs}-\tfrac{\bar g_Z^2}{M_Z^2}[Z_{\nu_L}]_{pr}[Z_{e_L}]^*_{st}\label{eq:LVLLRGE}\\[5pt]
    L^{V,LR}_{\nu e,prst}&= C_{le}^{prst} -\tfrac{\bar g_Z^2}{M_Z^2}[Z_{\nu_L}]_{pr}[Z_{e_R}]^*_{st}\label{eq:LVLRRGE}
\end{align}
For the $prst=ab12$ flavor combination under consideration we are therefore obtaining, while neglecting all terms quadratic in couplings, the following matching conditions:
\begin{align}
    L^{V,LL}_{\nu e,ab12}=\, &C_{ll}^{ab12}+C_{ll}^{12ab}-\tfrac{\bar g_2^2}{2M_W^2}\left(\delta_{a2}\delta_{b1}+\delta_{a2}C^{(3),b1}_{Hl}+\delta_{b1}C^{(3),a2}_{Hl}\right)\label{eq:LVLLmatching}\\[5pt]
    &+\tfrac{\bar g_Z^2}{M_Z^2}\left(\frac{1}{4}\delta_{ab}v_T^2\left[C_{Hl}^{(1),12*}+C_{Hl}^{(3),12*}\right]\right)\nonumber\\[5pt]
    L^{V,LR}_{\nu e,ab12}=\, &C_{le}^{ab12}+\tfrac{\bar g_Z^2}{M_Z^2}\left(\frac{1}{4}\delta_{ab}v_T^2C_{He}^{12*}\right)
\end{align}
We have therefore 5 SMEFT Operators at dimension six which the LEFT matches onto (full basis in \ref{tab:dim6SMEFTbasis}):
\begin{align}
    \op_{ll}^{prst} & =(\bar{l}_p \gamma^\mu l_r)(\bar{l}_s \gamma_\mu l_t) \label{Oll}\\[5pt]
    \op_{le}^{prst}&=(\bar{l}_p \gamma^\mu l_r)(\bar{e}_s \gamma_\mu e_t)\\[5pt]
    \op_{Hl}^{(1),pr}&=(H^\dagger i\overleftrightarrow{D}_\mu H)(\bar{l}_p \gamma^\mu l_r)\\[5pt]
    \op_{Hl}^{(3),pr} &=(H^\dagger i\overleftrightarrow{D}^I_\mu H)(\bar{l}_p \tau^I \gamma^\mu l_r)\\[5pt]
    \op_{He}^{pr}&=(H^\dagger i\overleftrightarrow{D}_\mu H)(\bar{e}_p \gamma^\mu e_r)
\end{align}
We see from the last term of Eq. \ref{eq:LVLLmatching} that the SMEFT operator $Q_{Hl}^{(1)}$ only produces LEFT operators with two identical neutrino generations (due to the $\delta_{ab}$) -- this is because the operator stems from a lepton number changing neutral current mediated by the $Z$ boson which couples the muon directly to the electron and decays itself into two neutrinos of the same generation.\par  
For the scalar and tensor operators, since they violate lepton number, the matching at dimension six is zero but at dimension seven we obtain from \parencite{liao2020extending}:
\begin{align}
L^{S,LL}_{\nu e,prst} &= -\frac{\sqrt{2}v}{8} \left( 2C_{\bar{e}lllH}^{stpr} + C_{\bar{e}lllH}^{sptr} + p \leftrightarrow r \right) \label{eq:SLLtoD7}\\
L^{S,LR}_{\nu e,prst} &= -\frac{\sqrt{2}v}{2} \left( C_{leHD}^{pt}\delta^{rs} + C_{leHD}^{rt}\delta^{ps} \right) \label{eq:SLRtoD7} \\
L^{T,LL}_{\nu e, prst} &= +\frac{\sqrt{2}v}{32} \left( C_{\bar{e}lllH}^{sptr} - C_{\bar{e}lllH}^{srtp} \right) \label{eq:TLLtoD7}
\end{align}
Notice that in this case the $v$ is the bare SM VEV -- corrections in terms of $v_T$ would start at order $1/\Lambda^5$ and are therefore neglected in this tree-level result. As we see in Eq.~\ref{eq:SLLtoD7} to Eq.~\ref{eq:TLLtoD7} there are two relevant operators at dimension 7 (full basis in \ref{tab:dim7SMEFTbl02} and \ref{tab:dim7SMEFTbl1-1}):
\begin{align}
    \mathcal{O}_{leHD}^{pr} &= \epsilon_{ij}\epsilon_{mn} (\overline{l_p^{iC}} \gamma_\mu e_r) H^j H^m iD^\mu H^n \label{eq:leHD}\\
    \mathcal{O}_{\bar{e}lllH}^{prst} &= \epsilon_{ij}\epsilon_{mn} (\overline{e_p} l_r^i) (\overline{l_s^{jC}} l_t^m) H^n \label{eq:elllH}
\end{align}
The specific matching for the seven operators that we derived $95\%$ CL intervals for (Tab.~\ref{tab:leftbounds}) is:
\begin{align}
        L^{V,LL}_{\nu e,2112} &=-\frac{2}{v_T^2} + C_{ll}^{2112} + C_{ll}^{1221}-2C^{(3),22}_{Hl}-2C^{(3),11}_{Hl}\label{eq:specMatchVLL2112}\\
        L^{V,LR}_{\nu e,2112} &= C_{le}^{2112}\label{eq:specMatchVLR2112}\\
        L^{V,LL}_{\nu e,3312} &= C_{ll}^{3312}+C_{ll}^{1233}+\frac{1}{4}v_T^2\tfrac{\bar g_Z^2}{M_Z^2}\left(C_{Hl}^{(1),12*}+C_{Hl}^{(3),12*}\right)\\
        L^{S,LL}_{\nu e,1112} &= -\frac{\sqrt{2}v}{4}\left(2C_{\bar elllH}^{1211}+C_{\bar e lllH}^{1121}\right)\\
        L^{S,LR}_{\nu e,1112} &=-\sqrt{2}vC_{LeDH}^{12}\\
        L^{S,LR}_{\nu e,1212} &=\frac{\sqrt{2}v}{2}C_{LeDH}^{22}\\
        L^{T,LL}_{\nu e,1212} &=\frac{\sqrt{2}v}{32}\left(C_{\bar elllH}^{1122}-C_{\bar elllH}^{1221}\right)\label{eq:specMatchTLL1221}
\end{align}
Let us analyze the very first matching conditions in Eq. \ref{eq:specMatchVLL2112} for $ L^{V,LL}_{\nu e,2112}$: The first term $-\tfrac{2}{v_T^2}$ amounts to the pure SM $W$-boson exchange -- however, in the spontaneously broken SMEFT the Higgs potential is modified by the operator $\mathcal{O}_H = (H^\dagger H)^3$, see Eq. \ref{eq:vtgaugecorr}. Since we are working in the  $\{\hat{m}_W, \hat{m}_Z, \hat{\alpha}\}$ input scheme, $v_T$ is not determined by the Higgs potential but by the physical mass of the $W$ boson:
\begin{equation}
    m_W^2 = \frac{1}{4} \bar{g}_2^2 v_T^2
\end{equation}
We can not really distinguish whether our LEFT deviation is due to an SM gauge shift, a vertex correction or the four-fermion operator. We  neglect the corrections to $\bar g$ that stem from the SMEFT in the broken phase and refer to the technical report of the LHC EFT working group in~\cite{brivio2022truncation}: Generally products of two dimension-six coefficients, including electroweak input-parameter shifts multiplying an already dimension-six contribution, are consequently of $\mathcal O(\Lambda^{-4})$ and can be neglected. 
For operators that do not interfere with the Standard-Model amplitude, the squared EFT amplitude is retained as the leading non-vanishing contribution. Such squared terms are separately gauge invariant and basis independent, but constitute only a partial contribution at the
corresponding higher order in the EFT expansion. Moreover, SM gauge shifts and vertex corrections are heavily restricted by Electroweak Precision Observables (EWPO) and LEP data -- we must unfortunately refrain from a detailed discussion at this point and only note that exclusively constraining the four-fermion operator is legitimate.\par
Since furthermore $C_{ll}^{2112}=C_{ll}^{1221*}$ we can simplify the matching to:
\begin{equation}
    \text{Re}[L^{V,LL}_{\nu e,2112}]\approx 2\text{Re}[C_{ll}^{2112}]
\end{equation}
The matching for $L^{V,LR}_{\nu e,2112} = C_{le}^{2112}$ \ref{eq:specMatchVLR2112} is much cleaner -- and moreover we have constraints for both its real and imaginary part.\par 
The matching of \(L_{\nu e,1112}^{S,LL}\) requires some
care. Here, the single-coefficient assumption at the LEFT level means
that only \(L_{\nu e,1112}^{S,LL}\) is varied -- it does not imply that the matching selects a single independent SMEFT flavour coefficient. The first index of \(\mathcal O_{\bar elllH}^{prst}\) labels the
charged-lepton singlet, whereas its last three indices label the three lepton doublets and obey non-trivial permutation relations. We therefore rewrite the raw flavour components in the irreducible flavour basis used by \texttt{D7RGESolver} \parencite{liao2025rge}. For the flavour multiset \((r,s,t)=(1,1,2)\), the totally antisymmetric component vanishes, leaving one symmetric and one mixed-symmetry coefficient. We define:
\begin{equation}
     C_S\equiv C_{\bar elllH}^{(S),1112}
 \qquad
 C_M\equiv C_{\bar elllH}^{(M),1112}
\end{equation}

The projector relations then imply:
\begin{equation}
     C_{\bar elllH}^{1211}=C_S-C_M
 \qquad
 C_{\bar elllH}^{1121}=C_S
 \qquad
 C_{\bar elllH}^{1112}=C_S+C_M
\end{equation}

It follows that:
\begin{align}
 L_{\nu e,1112}^{S,LL}
 &=-\frac{\sqrt2v}{4}
   \left(2C_{\bar elllH}^{1211}
         +C_{\bar elllH}^{1121}\right)
 \nonumber\\
 &=-\frac{\sqrt2v}{4}\left(3C_S-2C_M\right).
 \label{eq:SLL-SM-decomposition}
\end{align}
The weak-isospin contractions by the Levi-Civita tensors should not be
confused with this flavour decomposition: The latter follows from
permutations of the three lepton-doublet flavour indices.
\section{RGEs for the SMEFT}
In order to consistently constraint UV fields we have to bridge the energy scale upwards to them. From the electroweak scale at about $246\text{ GeV}$ to a few TeV, as estimated in Tab.~\ref{tab:leftbounds}, it is only about one order of magnitude, so we don't expect large corrections -- but on a conceptual level the running and mixing of operators is a key element of exploring energy scales and translating between them: UV fields might produce effective operators that look quite different at a low-energy process like muon decay.\par 
For the dimension-six operators, we are using the \texttt{wilson}~2.5.2 package \parencite{aebischer2018python} and a continuously fixed Warsaw weak basis, for the dim-7 operators the \texttt{D7RGESolver} from \parencite{liao2025rge} in the down basis.
The full RGEs are given in section \ref{sec:apprge} in the appendix. In Tab.~\ref{tab:smeft-rge-primary} we give the complete complete one-coefficient translation of the muon decay LEFT bounds into SMEFT bounds and subsequently into UV bounds together with the total relative change of the components between $\mu_\text{EW}$ and $\Lambda_\text{est}$. In Tab.~\ref{tab:smeft-rge-mixing-d7} we collect the mixing of the of the dimension-seven coefficient into the strongest generated coefficient and the Weinberg operator in order to check whether those values are consistent with current neutrino mass bounds. In Tab.~\ref{tab:smeft-rge-mixing-d6} we collect the mixing of the dimension-six coefficients into several of the coefficients they produce.\par 
Equation~\ref{eq:SLL-SM-decomposition} shows that the LEFT measurement
constrains one linear combination of two independent SMEFT flavour
coefficients:
\begin{equation}
     v^2L_{\nu e,1112}^{S,LL}
    =-\frac{\sqrt2}{4}
   \left(3z_{S,1112}-2z_{M,1112}\right)
    \qquad z_{S,M}=v^3C_{S,M}
\end{equation}
If both coefficients are allowed simultaneously, only this combination
is bounded. The two rows in Tab.~\ref{tab:smeft-rge-primary} instead
represent the mutually exclusive one-coefficient hypotheses
\(z_M=0\) and \(z_S=0\) -- they must not be interpreted as a simultaneous
fit. \par 
The naive contributions in Tab.~\ref{tab:induced-neutrino-masses}
range from $\mathcal{O}(10^2\,\mathrm{eV})$ to
$\mathcal{O}(10^4\,\mathrm{eV})$ and are consequently many orders of
magnitude above the sub-eV neutrino-mass scale indicated by oscillation
data and constrained by direct and cosmological measurements
\parencite{particle2024review}. The strongest direct measurement comes from the KATRIN experiment from the kinematics of tritium $\beta$-decay. From the combined analysis of its first five measurement campaigns, KATRIN obtains \cite{katrin2025direct}:
\begin{equation}
    m_\beta^2
    \equiv \sum_{i=1}^{3}\lvert U_{ei}\rvert^2m_i^2
    =-0.14^{+0.13}_{-0.15}\,\mathrm{eV}^2
    \quad\implies\quad
    m_\beta<0.45\,\mathrm{eV}
    \quad (90\%~\mathrm{CL})
\end{equation}

If one imposed the genuinely
single-operator boundary condition $C_5(\Lambda)=0$ at the UV scale,
the downward RGE evolution would therefore generate neutrino masses
that are incompatible with observation when the dimension-seven
coefficients saturate their muon-decay bounds. In such a scenario,
neutrino-mass constraints would be considerably stronger than those
obtained from the Michel parameters, unless the mixing is suppressed by
a symmetry or cancelled by an independent contribution to the Weinberg
operator. The values in the table should nevertheless not be interpreted
as direct low-energy predictions of the present calculation: The RGE
evolution was performed upwards with $C_5(\mu_{\mathrm{EW}})=0$, so that
the generated $C_5(\Lambda)$ instead quantifies the UV coefficient
required to compensate the mixing and recover the chosen electroweak
boundary condition \parencite{liao2025rge}.
% Auto-generated transcription of smeft_rge_results.md.
% The tables use only packages already loaded in main.tex:
% booktabs, graphicx, amsmath, amssymb.
%
% Definitions:
%   x_LEFT = v^2 L,  z_d = v^(d-4) C_d,  v = mu_EW = 246.22 GeV.
% All intervals are single-coefficient 95% CL intervals.

\begin{table}[htbp]
    \vspace*{1cm}
    \centering
    \resizebox{\textwidth}{!}{%
    \renewcommand{\arraystretch}{1.25}
    \begin{tabular}{lllllllll}
        \toprule
        LEFT
        & SMEFT
        & Matching
        & \(x_{\mathrm{LEFT}}(m_\mu)\)
        & \(x_{\mathrm{LEFT}}(\mu_{\mathrm{EW}})\)
        & \(z_d(\mu_{\mathrm{EW}})\)
        & \(\Lambda_{\mathrm{est}}\,[\mathrm{TeV}]\)
        & \(z_d(\Lambda_{\mathrm{est}})\)
        & \(\Delta_{\mathrm{RGE}}\)
        \\
        \midrule
        \multicolumn{9}{c}{dimension-six SMEFT coefficients}\\
        \midrule
        \(\operatorname{Re}L_{\nu e,2112}^{V,LL}\)
        & \(\operatorname{Re}C_{ll}^{2112}\)
        & \(x_{\mathrm{LEFT}}=2z_6\)
        & \([-0.0044,+0.0024]\)
        & \([-0.0044,+0.0024]\)
        & \([-0.0022,+0.0012]\)
        & \(5.277\)
        & \([-0.0022,+0.0012]\)
        & \(-0.41\%\)
        \\
        \(\operatorname{Re}L_{\nu e,2112}^{V,LL}\)
        & \(\operatorname{Re}C_{Hl}^{(3),11}\)
        & \(x_{\mathrm{LEFT}}=-2z_6\)
        & \([-0.0044,+0.0024]\)
        & \([-0.0044,+0.0024]\)
        & \([-0.0012,+0.0022]\)
        & \(5.277\)
        & \([-0.0012,+0.0023]\)
        & \(+4.84\%\)
        \\
        \(\operatorname{Re}L_{\nu e,2112}^{V,LL}\)
        & \(\operatorname{Re}C_{Hl}^{(3),22}\)
        & \(x_{\mathrm{LEFT}}=-2z_6\)
        & \([-0.0044,+0.0024]\)
        & \([-0.0044,+0.0024]\)
        & \([-0.0012,+0.0022]\)
        & \(5.277\)
        & \([-0.0012,+0.0023]\)
        & \(+4.84\%\)
        \\
        \(\operatorname{Re}L_{\nu e,2112}^{V,LR}\)
        & \(\operatorname{Re}C_{le}^{2112}\)
        & \(x_{\mathrm{LEFT}}=z_6\)
        & \([-0.0213,+0.0177]\)
        & \([-0.0213,+0.0177]\)
        & \([-0.0213,+0.0177]\)
        & \(1.687\)
        & \([-0.0211,+0.0176]\)
        & \(-0.95\%\)
        \\
        \(\operatorname{Im}L_{\nu e,2112}^{V,LR}\)
        & \(\operatorname{Im}C_{le}^{2112}\)
        & \(x_{\mathrm{LEFT}}=z_6\)
        & \([-0.0165,+0.0226]\)
        & \([-0.0165,+0.0226]\)
        & \([-0.0165,+0.0226]\)
        & \(1.638\)
        & \([-0.0164,+0.0224]\)
        & \(-0.94\%\)
        \\
        \(\operatorname{Re}L_{\nu e,3312}^{V,LL}\)
        & \(\operatorname{Re}C_{ll}^{3312}\)
        & \(x_{\mathrm{LEFT}}=2z_6\)
        & \([-0.1320,+0.1320]\)
        & \([-0.1320,+0.1320]\)
        & \([-0.0660,+0.0660]\)
        & \(0.958\)
        & \([-0.0656,+0.0656]\)
        & \(-0.63\%\)
        \\
        \(\operatorname{Re}L_{\nu e,3312}^{V,LL}\)
        & \(\operatorname{Re}C_{Hl}^{(1),12}\)
        & \(x_{\mathrm{LEFT}}=\operatorname{Re}z_6^\ast\)
        & \([-0.1320,+0.1320]\)
        & \([-0.1320,+0.1320]\)
        & \([-0.1320,+0.1320]\)
        & \(0.678\)
        & \([-0.1363,+0.1363]\)
        & \(+3.21\%\)
        \\
        \(\operatorname{Re}L_{\nu e,3312}^{V,LL}\)
        & \(\operatorname{Re}C_{Hl}^{(3),12}\)
        & \(x_{\mathrm{LEFT}}=\operatorname{Re}z_6^\ast\)
        & \([-0.1320,+0.1320]\)
        & \([-0.1320,+0.1320]\)
        & \([-0.1320,+0.1320]\)
        & \(0.678\)
        & \([-0.1342,+0.1342]\)
        & \(+1.63\%\)
        \\
        \midrule
        \multicolumn{9}{c}{dimension-seven SMEFT coefficients}\\
        \midrule
        \(\operatorname{Re}L_{\nu e,1112}^{S,LL}\)
        & \(\operatorname{Re}C_{\bar e LLLH}^{(S),1112}\)
        & \(x_{\mathrm{LEFT}}=-\frac{3\sqrt2}{8}z_7\)
        & \([+0.0364,+0.0689]\)
        & \([+0.0355,+0.0671]\)
        & \([-0.1265,-0.0669]\)
        & \(0.490\)
        & \([-0.1282,-0.0678]\)
        & \(+1.36\%\)
        \\
        \(\operatorname{Re}L_{\nu e,1112}^{S,LL}\)
        & \(\operatorname{Re}C_{\bar e LLLH}^{(M),1112}\)
        & \(x_{\mathrm{LEFT}}=\frac{\sqrt2}{4}z_7\)
        & \([+0.0364,+0.0689]\)
        & \([+0.0355,+0.0671]\)
        & \([+0.1003,+0.1898]\)
        & \(0.428\)
        & \([+0.1005,+0.1902]\)
        & \(+0.22\%\)
        \\
        \(\operatorname{Re}L_{\nu e,1112}^{S,LR}\)
        & \(\operatorname{Re}C_{LeDH}^{12}\)
        & \(x_{\mathrm{LEFT}}=-\sqrt2\,z_7\)
        & \([-0.0251,+0.0251]\)
        & \([-0.0244,+0.0244]\)
        & \([-0.0173,+0.0173]\)
        & \(0.953\)
        & \([-0.0184,+0.0184]\)
        & \(+6.40\%\)
        \\
        \(\operatorname{Re}L_{\nu e,1212}^{S,LR}\)
        & \(\operatorname{Re}C_{LeDH}^{22}\)
        & \(x_{\mathrm{LEFT}}=\frac{\sqrt2}{2}z_7\)
        & \([-0.0355,+0.0355]\)
        & \([-0.0345,+0.0345]\)
        & \([-0.0488,+0.0488]\)
        & \(0.674\)
        & \([-0.0511,+0.0511]\)
        & \(+4.81\%\)
        \\
        \(\operatorname{Re}L_{\nu e,1212}^{T,LL}\)
        & \(\operatorname{Re}C_{\bar e LLLH}^{(M),1122}\)
        & \(x_{\mathrm{LEFT}}=\frac{\sqrt2}{32}z_7\)
        & \([-0.0076,+0.0076]\)
        & \([-0.0077,+0.0077]\)
        & \([-0.1743,+0.1743]\)
        & \(0.441\)
        & \([-0.1747,+0.1747]\)
        & \(+0.23\%\)
        \\
        \bottomrule
    \end{tabular}%
    }
    \caption{%
        Complete one-coefficient translation of the muon decay LEFT bounds
        into SMEFT constraints and their subsequent evolution to the estimated
        UV scale. We define \(x_{\mathrm{LEFT}}=v^2L\) and
        \(z_d=v^{d-4}C_d\), with
        \(v=\mu_{\mathrm{EW}}=246.22\,\mathrm{GeV}\). The scale estimate assumes
        \(\lvert C_d\rvert=\Lambda^{4-d}\), so that
        \(\Lambda_{\mathrm{est}}
        =v/\max\lvert z_d\rvert^{1/(d-4)}\).
        The column \(\Delta_{\mathrm{RGE}}\) gives the relative change of the
        indicated real or imaginary component between
        \(\mu_{\mathrm{EW}}\) and \(\Lambda_{\mathrm{est}}\).
    }
    \label{tab:smeft-rge-primary}
\end{table}



\begin{table}[htbp]
    \centering
    \renewcommand{\arraystretch}{1.25}
    \begin{tabular}{lllll}
        \toprule
        EW-scale seed
        & \(\Lambda_{\mathrm{est}}\)
        & Change
        & D7 Cf.
        & Induced Weinberg Cf.
        \\
        \midrule
        \(\operatorname{Re}C_{\bar eLLLH}^{(S),1112}\)
        & \(0.490\,\mathrm{TeV}\)
        & \(+1.36\%\)
        & \(C_{\bar Q uLLH}^{3321}\;(3.55\times10^{-6}\%)\)
        & \(C_{LH5}^{12}\;(1.172\times10^{-8})\)
        \\
        \(\operatorname{Re}C_{\bar eLLLH}^{(M),1112}\)
        & \(0.428\,\mathrm{TeV}\)
        & \(+0.22\%\)
        & \(C_{\bar Q uLLH}^{3312}\;(3.76\times10^{-6}\%)\)
        & \(C_{LH5}^{12}\;(3.101\times10^{-9})\)
        \\
        \(\operatorname{Re}C_{LeDH}^{12}\)
        & \(0.953\,\mathrm{TeV}\)
        & \(+6.40\%\)
        & \(C_{\bar d LueH}^{3132}\;(0.0703\%)\)
        & \(C_{LH5}^{12}\;(1.610\times10^{-6})\)
        \\
        \(\operatorname{Re}C_{LeDH}^{22}\)
        & \(0.674\,\mathrm{TeV}\)
        & \(+4.81\%\)
        & \(C_{\bar d LueH}^{3232}\;(0.0527\%)\)
        & \(C_{LH5}^{22}\;(2.363\times10^{-6})\)
        \\
        \(\operatorname{Re}C_{\bar eLLLH}^{(M),1122}\)
        & \(0.441\,\mathrm{TeV}\)
        & \(+0.23\%\)
        & \(C_{\bar Q uLLH}^{3322}\;(1.99\times10^{-6}\%)\)
        & \(C_{LH5}^{22}\;(6.527\times10^{-9})\)
        \\
        \bottomrule
    \end{tabular}%
    \caption{%
        Complete dimension-seven mixing generated by
        \texttt{D7RGESolver} using integrated one-loop running in the down
        basis. The forth column displays the strongest induced dimension-seven operator and the percentages refers to the EWSB-normalized
        induced-to-input ratios. The final column separately displays the
        largest generated Weinberg coefficient as
        \(v\lvert C_5\rvert/(v^3\lvert C_{7,\mathrm{input}}\rvert)\);
        it is a dimensionless ratio rather than a percentage.
    }
    \label{tab:smeft-rge-mixing-d7}
\end{table}

\begin{table}[htbp]
    \centering
    \small
    \renewcommand{\arraystretch}{1.25}
    \begin{tabular}{lllllll}
        \toprule
        EW seed
        & \(\Lambda_{\mathrm{est}}\)
        & Notable?
        & Induced &$\tfrac{C_i(\Lambda_\text{est})}{C_\text{input}(\mu_\text{EW})}$ &Induced&$\tfrac{C_i(\Lambda_\text{est})}{C_\text{input}(\mu_\text{EW})}$
        \\
        \midrule
        \(\operatorname{Re}C_{ll}^{2112}\)
        & \(5.277\,\mathrm{TeV}\)
        & yes
        & \(C_{ll}^{1122}\)&\(4.33\%\)&
          \(C_{ll}^{2332}\)&\(0.536\%\)\\&&&
          \(C_{ll}^{1331}\)&\(0.536\%\)&
          \(C_{ll}^{2222}\)&\(0.368\%\)\\&&&
          \(C_{ll}^{1111}\)&\(0.368\%\)&
          \(C_{Hl}^{(3),22}\)&\(0.266\%\)\\&&&
          \(C_{Hl}^{(3),11}\)&\(0.266\%\)&
          \(C_{lq}^{(3),2211}\)&\(0.266\%\)
        \\\midrule
        \(\operatorname{Re}C_{Hl}^{(3),11}\)
        & \(5.277\,\mathrm{TeV}\)
        & yes
        & \(C_{H\Box}\)&\(1.75\%\)&
          \(C_{lq}^{(3),1133}\)&\(1.18\%\)\\&&&
          \(C_{uH}^{33}\)&\(0.953\%\)&
          \(C_{Hl}^{(3),22}\)&\(0.569\%\)\\&&&
          \(C_{Hl}^{(3),33}\)&\(0.569\%\)&
          \(C_{Hq}^{(3),11}\)&\(0.569\%\)\\&&&
          \(C_{Hq}^{(3),22}\)&\(0.569\%\)&
          \(C_{Hq}^{(3),33}\)&\(0.567\%\)
        \\\midrule
        \(\operatorname{Re}C_{Hl}^{(3),22}\)
        & \(5.277\,\mathrm{TeV}\)
        & yes
        & \(C_{H\Box}\)&\(1.75\%\)&
          \(C_{lq}^{(3),2233}\)&\(1.18\%\)\\&&&
          \(C_{uH}^{33}\)&\(0.953\%\)&
          \(C_{Hl}^{(3),11}\)&\(0.569\%\)\\&&&
          \(C_{Hl}^{(3),33}\)&\(0.569\%\)&
          \(C_{Hq}^{(3),11}\)&\(0.569\%\)\\&&&
          \(C_{Hq}^{(3),22}\)&\(0.569\%\)&
          \(C_{Hq}^{(3),33}\)&\(0.567\%\)
        \\\midrule
        \(\operatorname{Re}C_{le}^{2112}\)
        & \(1.687\,\mathrm{TeV}\)
        & no
        & \(C_{lequ}^{(1),1133}\)&\(0.00257\%\)&&
        \\\midrule
        \(\operatorname{Im}C_{le}^{2112}\)
        & \(1.638\,\mathrm{TeV}\)
        & no
        & \(C_{lequ}^{(1),1133}\)&\(0.00253\%\)&&
        \\\midrule
        \(\operatorname{Re}C_{ll}^{3312}\)
        & \(0.958\,\mathrm{TeV}\)
        & yes
        & \(C_{ll}^{1332}\)&\(2.12\%\)&
          \(C_{le}^{1222}\)&\(0.150\%\)\\&&&
          \(C_{le}^{1211}\)&\(0.150\%\)&
          \(C_{le}^{1233}\)&\(0.150\%\)\\&&&
          \(C_{lu}^{1211}\)&\(0.101\%\)&
          \(C_{lu}^{1222}\)&\(0.101\%\)\\&&&
          \(C_{lu}^{1233}\)&\(0.101\%\)&
          \(C_{ll}^{1222}\)&\(0.0774\%\)
        \\\midrule
        \(\operatorname{Re}C_{Hl}^{(1),12}\)
        & \(0.678\,\mathrm{TeV}\)
        & yes
        & \(C_{lu}^{1233}\)&\(1.03\%\)&
          \(C_{lq}^{(1),1233}\)&\(0.539\%\)\\&&&
          \(C_{le}^{1211}\)&\(0.0558\%\)&
          \(C_{le}^{1222}\)&\(0.0558\%\)\\&&&
          \(C_{le}^{1233}\)&\(0.0556\%\)&
          \(C_{lu}^{1211}\)&\(0.0373\%\)\\&&&
          \(C_{lu}^{1222}\)&\(0.0373\%\)&
          \(C_{ll}^{1222}\)&\(0.0281\%\)
        \\\midrule
        \(\operatorname{Re}C_{Hl}^{(3),12}\)
        & \(0.678\,\mathrm{TeV}\)
        & no
        & \(C_{lq}^{(3),1233}\)&\(0.431\%\)&
          \(C_{ll}^{1332}\)&\(0.176\%\)\\&&&
          \(C_{ll}^{1112}\)&\(0.0892\%\)&
          \(C_{ll}^{1222}\)&\(0.0892\%\)\\&&&
          \(C_{lq}^{(3),1211}\)&\(0.0878\%\)&
          \(C_{lq}^{(3),1222}\)&\(0.0870\%\)\\&&&
          \(C_{ll}^{1233}\)&\(0.0868\%\)&
          \(C_{lq}^{(3),1232}\)&\(0.0217\%\)
        \\
        \bottomrule
    \end{tabular}
        \caption{%
        Complete dimension-six operator mixing generated between the EW scale 
        \(\mu_{\mathrm{EW}}=246.22\,\mathrm{GeV}\) and the
        one-coefficient scale estimate \(\Lambda_{\mathrm{est}}\).
        Every percentage is the magnitude of the induced coefficient divided
        by the input coefficient. The numerical evolution uses
        \texttt{wilson}~2.5.2 in a continuously fixed Warsaw weak basis.
        ``Notable'' denotes mixing of at least \(1\%\); coefficients down to
        \(0.001\%\) are listed.
    }
    \label{tab:smeft-rge-mixing-d6}
\end{table}

\begin{table}[htbp]
    \centering
    \renewcommand{\arraystretch}{1.25}
    \begin{tabular}{llll}
        \toprule
        D7 Cf.
        & Weinberg
        & Neutrino mass-matrix entry
        & Naïve induced mass scale
        \\
        \midrule
        $C_{\bar e LLLH}^{(S),1112}$
        & $C_5^{12}$
        & $\lvert(M_\nu)_{e\mu}\rvert$
        & $193\text{--}365\,\mathrm{eV}$
        \\
        $C_{\bar e LLLH}^{(M),1112}$
        & $C_5^{12}$
        & $\lvert(M_\nu)_{e\mu}\rvert$
        & $76.6\text{--}145\,\mathrm{eV}$
        \\
        $C_{LeHD}^{12}$
        & $C_5^{12}$
        & $\lvert(M_\nu)_{e\mu}\rvert$
        & $\lesssim 6.86\,\mathrm{keV}$
        \\
        $C_{LeHD}^{22}$
        & $C_5^{22}$
        & $\lvert(M_\nu)_{\mu\mu}\rvert$
        & $\lesssim 28.4\,\mathrm{keV}$
        \\
        $C_{\bar e LLLH}^{(M),1122}$
        & $C_5^{22}$
        & $\lvert(M_\nu)_{\mu\mu}\rvert$
        & $\lesssim 280\,\mathrm{eV}$
        \\
        \bottomrule
    \end{tabular}
    \caption{%
        Naïve Majorana-neutrino mass scales associated with the
        Weinberg coefficients induced by the dimension-seven SMEFT
        operators. The values are obtained from
        $\lvert(M_\nu)_{pr}\rvert=v^2\lvert C_5^{pr}\rvert$
        using the endpoints of the corresponding $95\%$ confidence
        intervals from muon decay. The first two rows are ranges because the corresponding
        Michel-parameter intervals do not contain zero. For the remaining
        rows, whose intervals do contain zero, only the maximal mass scale
        at the endpoint of the interval is quoted. The coefficients
        $C_5^{12}$ and $C_5^{22}$ describe entries of the Majorana mass
        matrix in the charged-lepton flavor basis and should not be
        identified with individual mass eigenvalues $m_1$, $m_2$, or $m_3$.
    }
    \label{tab:induced-neutrino-masses}
\end{table}


